% Appendix: the model answers to Paper A, from exam/paper_a_solutions.md.

% The answers' headings are labelled ansa:q:N.
\renewcommand{\answerprefix}{ansa}

\chapter{Model Answers to Paper A}
\label{ansa:ch:answers}

Marks are shown per part. Method carries them: a correct derivation with a slip in it is worth more
than a correct answer with no working, and later parts consume earlier ones, so an error should be
followed through rather than penalised twice.

Where a part asks for VHDL, mark the \textbf{hardware described}, not the syntax. A listing with a
missing semicolon and the right circuit in it scores full marks; a listing that compiles and infers
a latch does not.

% ==========================================================================================
\answersection{1}{A table, a network, and a module}{12}

\answerpart{a}{4}
\textbf{Straight off the table.} Four rows are \code{1}, so four AND-terms, each carrying every
input directly where it is \code{1} in that row and primed where it is \code{0}:

\begin{console}
X = A'BC' + AB'C' + ABC' + ABC
\end{console}
\worth{1}

\textbf{Simplified.} The term \code{ABC'} is needed three times, which the idempotent law
\code{A + A = A} supplies for free:

\begin{textdrawing}
X = A'BC' + AB'C' + ABC' + ABC' + ABC' + ABC        (idempotent, twice)
  = (A' + A)BC' + (B' + B)AC' + AB(C' + C)          (distributive)
  = 1·BC' + 1·AC' + AB·1                            (complement)
  = BC' + AC' + AB                                  (identity)
\end{textdrawing}
\worth{2}

\textbf{The gate count.}

\begin{narrowtable}{@{}ll@{}}
  \thead{Form} & \thead{Gates} \\
  \midrule
  \code{A'BC' + AB'C' + ABC' + ABC} & 3 NOT, 4 three-input AND, 1 four-input OR = \textbf{8} \\
  \code{BC' + AC' + AB} & 1 NOT, 3 two-input AND, 1 three-input OR = \textbf{5} \\
\end{narrowtable}
\worth{1}

Worth noticing and worth a comment rather than a mark: \code{X = AB + AC' + BC'} is the
\textbf{majority vote} of \code{A}, \code{B} and \code{C'}. Check it against the table row by row
and every row agrees. That is the same observation \secref{c1:a2} makes when \code{A'B + AB'} turns
out to be one XOR gate: reading a table gives you a correct network, and recognizing the function
gives you a smaller one.

Accept any correct simplification path. Do not require these particular law names, but do require
\emph{some} justification per step: an answer that jumps from four terms to three with no working
is a remembered answer, not a derived one.

\answerpart{b}{3}
Both signals are active-low, so ``asserted'' means \code{0}:

\begin{textdrawing}
go = button_n' · enable_n'
\end{textdrawing}
\worth{1}

De Morgan, in the form \code{A'B' = (A + B)'}:

\begin{console}
go = (button_n + enable_n)'
\end{console}

\noindent which is a \textbf{NOR gate}, one gate for the whole function. \worth{1}

\begin{vhdlcode}
go <= (not button_n) and (not enable_n);   -- the expression as derived
go <= not (button_n or enable_n);          -- the same thing after De Morgan
\end{vhdlcode}

\noindent Synthesis produces \textbf{the same circuit for both}, and on the Cyclone~V that circuit
is not a NOR gate either: it is one two-input LUT loaded with the table
(Question~\ref{exa:q:5}(c)). The two lines are two spellings of one function, and the tool does not
care which you wrote. \worth{1}

This is the everyday use of De Morgan the book leans on: an active-low signal meeting logic written
for active-high, which in a book where every button and every reset is active-low is most of them.

\answerpart{c}{3}
\begin{vhdlcode}
library ieee;
use ieee.std_logic_1164.all;

entity gate_net is
    port(a, b, c: in  std_logic;
         x      : out std_logic);
end entity;

architecture behaviour of gate_net is
begin
    x <= (b and not c) or (a and not c) or (a and b);
end architecture;
\end{vhdlcode}

\noindent\textit{(2 marks: 1 for the entity with the ports in the stated order, 1 for an
architecture matching the candidate's own answer to (a))}

Naming the intermediate terms in internal signals is equally correct and synthesizes identically.

\textbf{What the clauses buy.} \code{std_logic} is not built into VHDL. It comes from the
\code{std_logic_1164} package in the \code{ieee} library, and so do the
\code{and}/\code{or}/\code{not} operators as they apply to it. Without those two lines the port
types do not exist and the file does not analyze. \worth{1}

Follow through an incorrect (a): an architecture that faithfully implements a wrong equation loses
nothing here.

\answerpart{d}{2}
\textbf{The concurrent assignment.} It describes a \textbf{wire}, or rather the gate driving one: a
standing, continuous relationship that holds at all times, exactly like the piece of hardware it
stands for. It is not executed once and it has no position in a sequence; \code{x} is
\code{a or b} for as long as the circuit exists. Two concurrent assignments in one architecture are
two pieces of hardware operating at the same time, not two steps. \worth{1}

\textbf{The two values.} \code{'U'} is uninitialized: nothing has driven this signal at all since
the simulation began. \code{'X'} is unknown: two things are driving it and they disagree. Neither is
a level a synthesized wire can carry, so neither is a value your design computed --- each is the
simulator telling you something is wrong, and the correct reaction is to find out what, not to work
around the symbol. \worth{1}

% ==========================================================================================
\answersection{2}{The map, the multiplexer, and the module used twice}{14}

\answerpart{a}{5}
The map, \code{AB} down the rows and \code{CD} across the columns, both in Gray-code order
\code{00, 01, 11, 10}:

\begin{narrowtable}{@{}ccccc@{}}
  \thead{\code{AB} \textbackslash{} \code{CD}} & \thead{\code{00}} & \thead{\code{01}} &
    \thead{\code{11}} & \thead{\code{10}} \\
  \midrule
  \code{00} & \textbf{1} &            &            & \textbf{1} \\
  \code{01} &            & \textbf{1} & \textbf{1} &            \\
  \code{11} &            & \textbf{1} & \textbf{1} &            \\
  \code{10} & \textbf{1} &            &            & \textbf{1} \\
\end{narrowtable}
\credit{1 mark for a correctly filled map in Gray-code order}

\textbf{Two groups of four.}

\begin{narrowtable}{@{}llll@{}}
  \thead{Group} & \thead{Cells} & \thead{Constant across the group} & \thead{Term} \\
  \midrule
  The four corners & \code{0000}, \code{0010}, \code{1000}, \code{1010} & \code{B = 0}, \code{D = 0}
    & \code{B'D'} \\
  The centre block & \code{0101}, \code{0111}, \code{1101}, \code{1111} & \code{B = 1}, \code{D = 1}
    & \code{BD} \\
\end{narrowtable}

\noindent The corner group is the one that has to be seen: it wraps \textbf{both} ways,
top-to-bottom and left-to-right, which is legal precisely because the Gray-code ordering makes the
last row adjacent to the first and the last column adjacent to the first. A candidate who instead
writes four groups of two here has read the map as a grid rather than as a torus, and has lost the
mark for this part but nothing later. \credit{2 marks: 1 per group, including the surviving
variables}

\begin{console}
X = B'D' + BD
\end{console}
\worth{1}

\textbf{The rest of the question.} Eight rows are \code{1}, so a direct Chapter~\ref{c1:ch}-style
reading gives \textbf{eight} AND-terms of four literals each. \code{A} and \code{C} do not appear
in the answer at all: the function does not depend on them. And \code{X = 1} exactly when \code{B}
and \code{D} are equal, so the whole thing is a single \textbf{XNOR} gate on \code{B} and \code{D}.
\worth{1}

That last step is the same one \secref{c1:a2} makes with \code{A'B + AB'}, and it is what the
technique is for: the map found a two-term equation, and recognizing the two-term equation found a
one-gate circuit.

\answerpart{b}{4}
\begin{vhdlcode}
architecture behaviour of mux4 is
begin
    MUX_PROCESS: process(inputs, sel) is
    begin
        case (sel) is
            when "00"   => x <= inputs(3);
            when "01"   => x <= inputs(2);
            when "10"   => x <= inputs(1);
            when "11"   => x <= inputs(0);
            when others => x <= '0';
        end case;
    end process;
end architecture;
\end{vhdlcode}

\noindent\textit{(2 marks: 1 for the four branches mapping the right way round, 1 for a
\code{process} with both \code{inputs} and \code{sel} in the sensitivity list. Deduct nothing for a
label, or for its absence.)}

\textbf{Why the \code{case} needs a process.} \code{case}, like \code{if} and \code{for}, is a
\textbf{sequential} statement, legal only inside a process. An architecture body holds concurrent
statements: standing descriptions of hardware, all active at once. A process is how a block of
sequential statements is packaged up so that it can be one of those concurrent statements.
\worth{1}

\textbf{What \code{when others} protects against.} It is a \textbf{language} requirement, not a
defensive one. VHDL insists a \code{case} cover every value of its selector's \emph{type}, and
\code{sel} is a \code{std_logic_vector}, whose elements have nine values each rather than two, so
\code{"00"}, \code{"01"}, \code{"10"} and \code{"11"} do not make the \code{case} exhaustive. It is
not a claim about the wire: a synthesized net carries a \code{0} or a \code{1} and nothing else, and
the other values belong to the simulator, where a signal sitting at \code{'U'} or \code{'X'} is a
bug rather than a level. \worth{1}

An answer arguing that \code{when others} ``catches an illegal selector on the board'' has the wrong
reason and should get no more than half the mark, even though the line it writes is identical.

\answerpart{c}{3}
\begin{vhdlcode}
architecture structure of mux4 is
signal x_high, x_low: std_logic;
begin
    -- sel(0) chooses within each pair.
    mux2_high: entity work.mux2
        port map(inputs(2), inputs(3), sel(0), x_high);

    mux2_low: entity work.mux2
        port map(inputs(0), inputs(1), sel(0), x_low);

    -- sel(1) chooses between the two pairs.
    mux2_out: entity work.mux2
        port map(x_low, x_high, sel(1), x);
end architecture;
\end{vhdlcode}

\noindent\textit{(2 marks: 1 for three instances with unique labels and two internal signals, 1 for
getting the \code{a}/\code{b} order right against \code{mux2}'s \code{sel = '1'} selecting
\code{a})}

Check it against the specification one selector value at a time, which is the check worth writing
down: \code{"00"} takes \code{x_high} and within it \code{inputs(3)}; \code{"01"} takes
\code{x_high} and \code{inputs(2)}; \code{"10"} takes \code{x_low} and \code{inputs(1)};
\code{"11"} takes \code{x_low} and \code{inputs(0)}.

\textbf{The rest.} The FPGA ends up with \textbf{three} copies of \code{mux2}'s logic, exactly as
writing the logic three times would. Instantiating a module more than once saves source, not
silicon.

And yes, this architecture and part (b)'s describe the same hardware, or near enough that the tool
will not distinguish them: both are a four-input multiplexer, and a fitter given either one packs it
into LUTs the same way. (b) describes the \textbf{intent} and lets the tool find the gates; (c)
describes the \textbf{structure} and hands the tool a wiring diagram. That gap is most of what a
hardware description language buys you. \worth{1}

\answerpart{d}{2}
The testbench connects to the module \textbf{by position}, not by name, so the entity's port
declaration order is part of its interface. A wrong order fails in one of two ways, and neither says
``wrong order'':
\begin{itemize}
  \item \textbf{Different types or widths.} Analysis stops with a message like
    \code{can't associate "sel" with port "sel"}, pointing at the testbench's \code{port map} line.
    You find out at \code{ghdl -a}, immediately.
  \item \textbf{Same type}, several \code{std_logic} ports swapped, say. It analyzes and elaborates
    quite happily, then fails at run time as an ordinary assertion failure that reads exactly like a
    logic bug in the architecture. You find out at \code{ghdl -r}, after spending time in the wrong
    file.
\end{itemize}
\credit{1 mark per case, including when it is caught}

The practical consequence, which is why the exercises spell their ports out in order: a testbench
that fails on the very first case it checks is a reason to reread the port order before hunting
through the architecture.

% ==========================================================================================
\answersection{3}{What a clock edge does, and when an assignment happens}{13}

\answerpart{a}{3}
\textbf{The latch} is \textbf{level-sensitive}: while \code{enable} is high it is transparent and
\code{Q} follows \code{D} continuously; while \code{enable} is low it holds. It changes its output at
any time during the entire enabled window.

\textbf{The flip-flop} is \textbf{edge-triggered}: it samples \code{D} at the single \emph{instant}
the clock transitions and holds that value for the rest of the period, whatever \code{D} does.
\worth{1}

\textbf{Why synchronous designs use the flip-flop.} Because the latch is transparent for a whole
clock phase, any glitch on \code{D} during that phase passes straight through, and guaranteeing data
is stable across an entire phase is far harder than guaranteeing it around one instant. Every
flip-flop changing at the same well-defined instants is also what makes a large circuit's timing
analyzable: you can reason one clock cycle at a time instead of tracking every signal continuously.
\worth{1}

\textbf{\code{D} rises and falls between two edges.} The latch follows it up and back down, so
\code{Q} shows the pulse (during whatever part of it the clock is high). The flip-flop shows
\textbf{nothing at all}: no edge happened while \code{D} was high, so \code{Q} never moves. Both are
behaving correctly; they are answering different questions. \worth{1}

\answerpart{b}{4}
A signal assignment does not take effect when its line runs. It \textbf{schedules} a value, applied
only once the process suspends, so every read during the pass returns the value the signal held
\emph{coming into} the edge. \code{c <= b} therefore reads the old \code{b}, not the one being
scheduled on the line below.

\begin{narrowtable}{@{}lcccccc@{}}
  \thead{Edge} & \thead{1} & \thead{2} & \thead{3} & \thead{4} & \thead{5} & \thead{6} \\
  \midrule
  \code{a} & 1 & 0 & 1 & 1 & 0 & 0 \\
  \code{b} after & \textbf{1} & \textbf{0} & \textbf{1} & \textbf{1} & \textbf{0} & \textbf{0} \\
  \code{c} after & \textbf{0} & \textbf{1} & \textbf{0} & \textbf{1} & \textbf{1} & \textbf{0} \\
\end{narrowtable}
\credit{2 marks for the table; 1 if \code{b} is right and \code{c} is shifted the wrong way or not
at all}

\textbf{The hardware.} Two flip-flops in a chain. Every signal assigned inside a clocked process
becomes a flip-flop, one bit each, and here \code{a} reaches \code{b} after one edge and \code{c}
after two. \worth{1}

\textbf{Swapping the lines changes nothing.} Neither assignment can observe the other's result,
because both read the pre-edge values and both take effect together when the process suspends. Order
does not matter between assignments to \emph{different} signals in one clocked process, which is
the opposite of the sequential reading the word ``process'' invites, and it is why a chain of
flip-flops can be written in any order. \worth{1}

Two related rules are worth a comment if a candidate volunteers them: assigning the \textbf{same}
signal twice in one pass keeps only the last assignment, and a \code{variable} behaves the other way
round on both counts, which is Question~\ref{exa:q:5}.

\answerpart{c}{3}
\begin{vhdlcode}
signal req_prev: std_logic;
...
    -- One gate: high for the single cycle where req is high and was not.
    req_edge <= req and not req_prev;

    REQ_PROCESS: process(clock, reset_n) is
    begin
        if (reset_n = '0') then
            req_prev <= '0';
        elsif (rising_edge(clock)) then
            req_prev <= req;
        end if;
    end process;
\end{vhdlcode}

\noindent\textit{(2 marks: 1 for the flip-flop holding the previous value with an asynchronous
active-low reset, 1 for \code{current AND NOT previous} as the rising-edge form. A registered
\code{req_edge} inside the process is also correct hardware, one cycle later; accept it if the
candidate says so.)}

\textbf{How long it is high.} Exactly \textbf{one clock cycle}: the one immediately after the
transition, because \code{req_prev} only catches up at the \emph{next} edge, at which point
\code{req} and \code{req_prev} agree again and the gate goes low.

\textbf{Why \code{'0'} here.} The reset value has to be the signal's \textbf{idle} level, so that
coming out of reset does not look like a transition. \code{req} is active-high and idles low, so
\code{req_prev} resets to \code{'0'}. \code{led_toggle}'s buttons are active-\textbf{low} and idle
at \code{1}, so its \code{button_prev_n} resets to \code{"11"}; reset it to \code{"00"} and the
circuit comes out of reset claiming both buttons are already held down. \worth{1}

\answerpart{d}{3}
\textbf{The rule.} \code{clock} is the only signal in the sensitivity list, unless the design has an
asynchronous reset, in which case that reset joins it and nothing else does.

\textbf{The reason.} The process has no business reacting to anything but the events that can
change its outputs. A flip-flop's output changes at a clock edge, and an asynchronously reset one
also changes the moment the reset is asserted, with no clock edge involved, so the reset has to be
able to wake the process on its own. Everything else the process reads is sampled \emph{at} the edge
and has no business waking it. \worth{1}

\textbf{Adding \code{d} to the list.} In \textbf{simulation}, the process now also runs whenever
\code{d} changes; since \code{rising_edge(clock)} is false on those runs, nothing is assigned, and
the behaviour is unchanged. In \textbf{hardware} nothing changes either: the sensitivity list is not
synthesized. So this one costs nothing but noise, and it is worth saying so plainly, because the
reflex is to assume every sensitivity-list mistake is fatal. \worth{1}

\textbf{Leaving \code{reset_n} out.} Now the two disagree. In \textbf{simulation} the process only
runs on \code{clock} events, so the reset is no longer asynchronous: asserting it does nothing until
the next clock edge, and a reset asserted and released between two edges is missed entirely. In
\textbf{hardware} the tool reads the
\code{if (reset_n = '0') then ... elsif (rising_edge(clock))} structure and builds the asynchronous
reset anyway, or refuses the process. Either way the simulator and the chip no longer agree, which
is worse than being wrong in one place. \worth{1}

% ==========================================================================================
\answersection{4}{The input you do not control}{13}

\answerpart{a}{3}
\textbf{Setup and hold} are the two halves of a window around the active clock edge: the
\textbf{setup time} is how long \code{D} must already be stable \emph{before} the edge, and the
\textbf{hold time} is how long it must stay stable \emph{after} it.

\textbf{What happens inside.} A flip-flop is built from cross-coupled gates that snap to \code{0} or
\code{1} when the clock samples. An input changing inside the window can leave that feedback loop
balanced between the two states, neither fully high nor fully low, for longer than expected:
\textbf{metastability}. A ball balanced on top of a hill rolls down eventually, but when it commits,
and to which side, is not predictable from the moment it was placed there, and while it is still
balanced its position is not a valid digital value at all. \worth{1}

\textbf{Why an asynchronous input cannot be promised to miss the window.} Because you do not control
when it changes. A push button, an external reset, or a line from another clock domain changes when
the outside world decides to, with no fixed relationship to your clock, so no amount of design makes
the window unreachable. \worth{1}

\textbf{What the actual failure is.} Not a wrong answer. Sampling a signal that genuinely was in the
middle of a transition and getting either level back is a legitimate outcome and is harmless: both
answers were true at some instant inside the window. The failure is the \textbf{delay before there
is an answer at all}, and specifically what happens if the unresolved voltage reaches downstream
logic before it settles, because then different gates can legitimately read it differently at the
same instant. One part of the circuit increments and another does not; a state machine lands in an
impossible state. The disagreement is the bug. \worth{1}

An answer that says ``it gives the wrong value'' has missed the whole point of \secref{c4:a2} and
scores nothing for this third mark, however fluent the rest is.

\answerpart{b}{4}
\begin{textdrawing}
async_in
   │
   ▼
 [FF1] ──────► [FF2] ──────► synchronized, safe to use
   (may go        (extremely likely
   metastable)     to be stable)
\end{textdrawing}

\noindent Only \textbf{FF2's} output may be used elsewhere. FF1 is the flip-flop exposed to the
asynchronous input, so it is the one that may go metastable; its output has an entire clock period
to settle before FF2 samples it, and by then it almost always has. FF1's output used directly is the
unresolved node itself, which is the failure in (a). \worth{1}

\textbf{The terms.}

\begin{booktable}{@{}p{4.5em}L@{}}
  \thead{Term} & \thead{What it is} \\
  \midrule
  \code{t_r} & The resolution time a stage is given, roughly one clock period \\
  \code{tau} & A property of the flip-flop: how fast a metastable value decays \\
  \code{T_w} & A property of the flip-flop: the width of the window around the edge in which an
    input can provoke metastability at all \\
  \code{f_clock} & The clock rate of the receiving domain \\
  \code{f_data} & How often the asynchronous input actually changes \\
\end{booktable}
\credit{2 marks, or 1 for naming three or four of the five}

\textbf{Which input makes the case for a third stage.} The \textbf{10\,MHz line}. It is
\code{f_data} that separates the two, and \code{1 Hz} against \code{10^7 Hz} is \textbf{seven orders
of magnitude} in that term alone; \code{t_r}, \code{tau}, \code{T_w} and \code{f_clock} are
identical for both, because they are properties of the flip-flop and the receiving clock, not of the
input. So the case for a third stage is a question about the input, not about the clock.

\textbf{What it costs.} One more clock cycle of latency, against a numerator that grows
\emph{exponentially} with each added \code{t_r}. That asymmetry is why extra stages exist at all.
\worth{1}

\answerpart{c}{3}
\begin{vhdlcode}
library ieee;
use ieee.std_logic_1164.all;

entity reset_sync is
    port(clock, reset_n: in  std_logic;
         reset_s2_n    : out std_logic);
end entity;

architecture behaviour of reset_sync is
signal reset_s1_n: std_logic;
begin
    process(clock, reset_n) is
    begin
        if (reset_n = '0') then
            reset_s1_n <= '0';
            reset_s2_n <= '0';
        elsif (rising_edge(clock)) then
            reset_s1_n <= '1';
            reset_s2_n <= reset_s1_n;
        end if;
    end process;
end architecture;
\end{vhdlcode}

\noindent\textit{(2 marks: 1 for both stages forced to \code{'0'} in the asynchronous branch, 1 for
the clocked branch shifting a constant \code{'1'} in rather than shifting \code{reset_n} in)}

Nothing here reads \code{reset_s2_n} back, so under VHDL-93 it can be driven straight from the
process, as \secref{c4:a5}'s listing does. A candidate who routes it through an extra internal
signal is equally correct, just carrying one declaration that buys nothing.

\textbf{The pattern: assert asynchronously, release synchronously.}
\begin{itemize}
  \item \textbf{Assert.} When \code{reset_n} drops, both stages are forced to \code{0} immediately,
    with no dependency on the clock. A reset that waited for a clock edge would defeat its own
    purpose: it has to take effect the instant it is needed, clock or no clock.
  \item \textbf{Release.} When \code{reset_n} returns to \code{1} the reset condition does not
    vanish everywhere at once. It shifts out through the two stages like any other asynchronous
    input, so the signal the design actually uses releases cleanly on a clock edge along with
    everything else.
\end{itemize}
\credit{1 mark for both halves}

\textbf{Without the synchronous release}, in a design of a dozen flip-flops --- worth a comment
rather than a mark, but the sentence the other two marks are really testing for: the reset's rising
edge arrives at each of them at a slightly different time relative to the clock, so some leave reset
on one edge and some on the next, and for one cycle the design holds a combination of state that no
correct execution ever produces. \code{serial_rx8} is the concrete case the book gives: its shift
register and its bit counter are a dozen flip-flops between them, and a receiver that starts
counting a byte one cycle before it starts shifting one has been broken by its own reset.

\answerpart{d}{3}
\textbf{What is safe.} Each bit individually. Every bit gets a full clock period to resolve, so no
\emph{single} bit is metastable when the receiving domain reads it. Metastability, per bit, is
solved.

\textbf{What is not.} The bits are not safe \textbf{together}. Each one independently resolves to
the old or the new value on the same edge, so the receiver can sample a combination that was never
on the bus. The double-flop answer synchronizes \textbf{one bit}, and that word is the whole limit.
\worth{1}

\textbf{A concrete value.} The counter stepping from \code{0111} to \code{1000} changes all four
bits at once. Each synchronizer independently lands on the old or the new bit, so the receiver can
read \code{1111}, or \code{0000}, or any of the other fourteen combinations --- values the counter
held neither before nor after the step. \worth{1}

\textbf{The three fixes.}
\begin{itemize}
  \item A \textbf{handshake}: the sender holds the bus still and raises a single-bit request, which
    \emph{is} safe to synchronize; the receiver reads the stable bus and acknowledges.
  \item A \textbf{FIFO} with synchronized pointers, which is the same idea industrialized.
  \item \textbf{Gray coding} the counter, so that only one bit ever changes per step and there is
    no combination to get wrong.
\end{itemize}
Gray code is the one that removes the problem \textbf{by construction}: there is no instant at which
two bits are in flight, so a per-bit synchronizer can only return the old value or the new one,
both of which are values the counter really held. The other two work by making sure nobody looks
while the bus is moving. \worth{1}

None of this is needed anywhere in this book, which never crosses a clock domain with a bus.
Knowing that the two-flop answer stops at one bit is the examinable part.

% ==========================================================================================
\answersection{5}{The accumulator that isn't, and the latch nobody asked for}{13}

\answerpart{a}{4}
Two rules do all the work: a signal assignment \textbf{schedules} rather than applies, so every read
during the pass returns the value held \emph{before the pass started}; and assigning the same signal
more than once in one pass keeps only the \textbf{last} assignment.

\code{bits'range} on \code{(7 downto 0)} runs high to low, and \code{bits = "01000000"}, so only
\code{bits(6)} is \code{'1'}.

\begin{narrowtable}{@{}lcccc@{}}
  \thead{Step} & \thead{\code{bits(i)}} & \thead{Read from \code{any_s}} & \thead{Scheduled} &
    \thead{Survives?} \\
  \midrule
  \code{any_s <= '0'} & -- & -- & \code{'0'} & no \\
  \code{i = 7} & \code{0} & \code{'0'} & \code{'0'} & no \\
  \code{i = 6} & \code{1} & \code{'0'} & \code{'1'} & no \\
  \code{i = 5} & \code{0} & \code{'0'} & \code{'0'} & no \\
  \code{i = 4} & \code{0} & \code{'0'} & \code{'0'} & no \\
  \code{i = 3} & \code{0} & \code{'0'} & \code{'0'} & no \\
  \code{i = 2} & \code{0} & \code{'0'} & \code{'0'} & no \\
  \code{i = 1} & \code{0} & \code{'0'} & \code{'0'} & no \\
  \code{i = 0} & \code{0} & \code{'0'} & \code{'0'} & \textbf{yes} \\
\end{narrowtable}

\noindent\textit{(2 marks for the trace, of which 1 is for reading \code{'0'} in \textbf{every} row
rather than the previous row's scheduled value)}

\code{found} ends up \code{'0'}. The correct answer is \code{'1'}: \code{bits(6)} is set. Nine
assignments were computed and eight of them were thrown away.

\textbf{Why this is not an accumulator.} Every iteration reads the same frozen value, so no
iteration can see the previous one's result, and only the textually last assignment survives. The
final value depends on exactly two things: whatever \code{any_s} already held, and \code{bits(0)}.
The loop is doing no accumulating at all --- it is computing the same expression eight times with a
different last operand and discarding all but one of the answers. \worth{1}

\textbf{When it happens to be right.} The architecture computes \code{old_any_s or bits(0)} where
the correct answer is the OR of all eight bits. So:
\begin{itemize}
  \item \textbf{Whenever \code{bits(0) = '1'} it is right}, unconditionally, because both
    expressions are then \code{'1'}. That is \textbf{half of all 256 inputs}.
  \item When \code{bits(0) = '0'} it is right only by luck: when the value it happened to be
    carrying already matches the correct answer.
\end{itemize}
That is what makes this bug dangerous rather than merely wrong. \code{"11111111"} passes.
\code{"00000001"} passes. Anything ending in a \code{1} passes. A smoke test of a few ``obvious''
inputs has a good chance of reporting success, and the design fails on the ones nobody typed.
\worth{1}

A candidate who adds that in a real simulation \code{any_s} starts at \code{'U'} rather than
\code{'0'}, so the first pass propagates \code{'U'} and Chapter~\ref{c1:ch}'s ``a signal sitting at
\code{'U'} is a bug, not a value'' fires immediately, has understood more than the question asked
and should be credited within the marks available.

So has one who points out that the surviving assignment is \code{any_s <= any_s or bits(0);} with
\code{any_s} absent from the sensitivity list, so what synthesis is being handed is a
\textbf{combinational feedback loop}: an output fed back to its own input through a gate, with no
flip-flop in the path. Quartus reports that rather than building it. ``Compiles cleanly'' here means
GHDL analyzes and simulates it, which is exactly the gap the question is about.

\answerpart{b}{3}
\begin{vhdlcode}
architecture behaviour of any_set is
signal any_s: std_logic;
begin
    found <= any_s;

    ANY_PROCESS: process(bits) is
        variable acc: std_logic;
    begin
        acc := '0';
        for i in bits'range loop
            acc := acc or bits(i);
        end loop;
        any_s <= acc;
    end process;
end architecture;
\end{vhdlcode}

\noindent\textit{(1 mark: a \code{variable}, assigned with \code{:=}, cleared at the top of every
activation, handed to the signal once at the end)}

\textbf{Clock cycles: zero.} This is a combinational process. The loop's bounds are fixed at
elaboration, so the synthesis tool \textbf{unrolls} it: the eight iterations become a chain of OR
gates that all settle within one propagation delay, in the same clock cycle. There is no counter for
\code{i}, no branch, and no iteration in time. The loop writes gates in three lines; it does not do
eight things in sequence. \worth{1}

\textbf{Storage: none.} \code{acc} reads like a running total but never has to survive a clock edge.
Each of its values is simply the wire between two of those OR gates, which is exactly why a
\code{variable} is the right tool: it never outlives one pass. \worth{1}

An answer of ``eight clock cycles'' is the C reading of the loop and scores nothing for the second
and third marks, however well the code is written.

\answerpart{c}{3}
A Boolean function becomes a \textbf{look-up table}: a tiny memory, loaded with the function's
truth table, which then computes it. Nothing becomes a gate.

The fabric is made of \textbf{LUTs and flip-flops}, one bit of clocked storage sitting beside each
LUT. On the DE0-CV's Cyclone~V they are packaged as an \textbf{ALM}, holding a fracturable six-input
LUT plus four flip-flops, and the ALM is what the Quartus fitter report counts --- which is why it
talks about ALMs rather than gates. \worth{2}

\textbf{What the map bought.} On this board, in LUT count, nothing. A four-input function costs one
LUT whether you minimized it or not, because in every case it is the same block holding a different
table. What Question~\ref{exa:q:2}'s map actually bought was the \emph{knowledge} that \code{A} and
\code{C} are irrelevant and the function is one XNOR --- which no amount of LUT packing would have
told you, and which is the sort of thing you need before you can tell whether a specification is the
one you meant.

What minimization still buys anywhere: it is how you understand a circuit; it is exactly right for
the discrete-logic and ASIC flows Karnaugh maps were invented for; and reducing a function below the
LUT's input width is what stops it needing a second level of LUTs. \worth{1}

\answerpart{d}{3}
\textbf{What the tool built: a latch}, one per bit of \code{leds}, exactly the D latch of
\secref{c3:a2}.

\textbf{Why it had no choice}, in two steps:
\begin{enumerate}
  \item VHDL requires a signal to keep its old value if nothing assigns a new one. The
    \code{when others} branch assigns nothing, so on that path the design has said \code{leds} must
    \textbf{remember} its previous value.
  \item Remembering is not something combinational logic can do. The only element that can is a
    latch, so that is what gets built.
\end{enumerate}
\worth{1}

\textbf{Two fixes.}

\begin{vhdlcode}
        when others => leds <= "0000";      -- assign on every path, in the branch
\end{vhdlcode}

\begin{vhdlcode}
    leds <= "0000";                          -- a default before the case; every branch overrides it
    case sel is
        ...
\end{vhdlcode}
\credit{1 mark for both}

The second is the one to prefer, and the reason is maintenance rather than hardware: it scales to a
\code{case} with twenty branches and stays correct when somebody adds a twenty-first, whereas the
first has to be got right again every time the \code{case} grows. Both synthesize to the same
combinational logic.

\textbf{Why a warning is more dangerous than an error.} An error stops you. A warning does not, and
the design will often appear to work on the board, because the latch happens to be holding the
right value most of the time --- so the bug ships, and it surfaces later as an intermittent fault in
a circuit that ``compiled cleanly''. ``It compiled with no errors'' promises a great deal less about
an FPGA design than the same sentence does about a C program. \worth{1}

% ==========================================================================================
\answersection{6}{Counting, shifting, and knowing when a byte has arrived}{12}

\answerpart{a}{4}
\textbf{After \code{15} it returns to \code{0}}, with nothing in the circuit saying to: no
comparison, no reset logic, no ``if the count is at maximum'' anywhere.

\textbf{The element that performs the wrap} is the adder's \textbf{carry-out}: the fifth bit. The
adder produced \code{10000}, the fifth bit had nowhere to go, and the four bits that survived are
\code{0000}. Discarding that bit \emph{is} the wraparound --- ``the count wraps because it runs out
of bits'' is not a figure of speech, and you can point at the bit it ran out of. \worth{2}

\textbf{What has to be true of the range.} The upper bound must be \code{2^N - 1}. \code{0 to 15}
wraps for free; \code{0 to 9} does not, because \code{10} is not a power-of-two boundary, so a
modulo-10 counter needs an explicit comparison against \code{9} and an explicit clear. \worth{1}

\textbf{Which is modelling the hardware.} The drawing and the chip. There is no sixteenth value for
four bits to take, so they land on \code{0000}. GHDL objects because \code{natural range 0 to 15} is
a \textbf{promise the designer made} and \code{16} breaks it; a range-constrained integer is a claim
about the value, not a description of a wire.

\textbf{The rule that follows} is not ``trust the hardware''. It is that anything you intend to
verify should say what it means: write the comparison out, the way \code{timer.vhd} and
\code{serial_rx8.vhd} do, and the drawing, the simulator and the chip all agree. \worth{1}

\answerpart{b}{3}
\code{shift_reg(6 downto 0)} is everything except the top bit, moved up one place; the top bit falls
off; \code{serial_in} fills the vacated bit 0.

\begin{narrowtable}{@{}ccc@{}}
  \thead{Edge} & \thead{\code{serial_in}} & \thead{\code{shift_reg} after} \\
  \midrule
  -- & -- & \code{00000000} \\
  1 & \code{1} & \code{00000001} \\
  2 & \code{0} & \code{00000010} \\
  3 & \code{1} & \code{00000101} \\
  4 & \code{1} & \code{00001011} \\
\end{narrowtable}
\credit{2 marks; 1 if the sequence is right but shifted the wrong way}

The first bit sent is now at \textbf{bit 3}, and after eight edges it will be at \textbf{bit 7}: the
first-arriving bit ends up in the most significant position, which is what ``most significant bit
first'' means for this book's receivers.

The mirror image is

\begin{vhdlcode}
shift_reg <= serial_in & shift_reg(7 downto 1);
\end{vhdlcode}

\noindent which moves every bit toward the \textbf{least} significant end, takes data in at bit 7
and drops it from bit 0, and so leaves the first-arriving bit in bit 0. That is the
least-significant-first order a real UART uses. It reads almost identically at a glance, which is
why the direction is worth checking before assuming. \worth{1}

\answerpart{c}{2}
\begin{narrowtable}{@{}lll@{}}
  & \thead{Data in} & \thead{Data out} \\
  \midrule
  \textbf{SIPO} & one bit per clock & all N bits at once \\
  \textbf{PISO} & all N bits at once, loaded & one bit per clock \\
\end{narrowtable}

\noindent\textbf{SIPO} to receive a stream of serial bits: it accumulates them and hands you the
whole word. \textbf{PISO} to drive a chain of LEDs from a fixed pattern: it takes the word and sends
it out one wire. \worth{1}

\textbf{What differs between the architectures}: almost nothing. The shift expression is
\textbf{identical} --- same direction, same line of VHDL. PISO adds a \code{load} branch to
distinguish loading a new parallel value from shifting the existing one out, and taps
\code{shift_reg(N-1)} as a serial output where SIPO exposes the whole vector. The two names describe
how you \emph{wire up} a shift register, not two kinds of hardware. \worth{1}

\answerpart{d}{3}
\textbf{Why a pulse.} The event being reported is instantaneous: a byte has just completed. A level
would say ``a byte completed at some point'', which a consumer cannot count. \code{data_ready} is
driven low unconditionally at the top of the clocked branch and raised only on the edge completing a
byte, which is the same single-cycle shape as Chapter~\ref{c3:ch}'s edge detector. \worth{1}

\textbf{What it requires of \code{data_out}.} That it \textbf{holds} its value afterwards. A consumer
that misses the pulse has missed the \emph{event}, not the data, and the byte is still there. A real
peripheral would add a holding register or a small FIFO here so a slow consumer cannot lose data at
all. \worth{1}

\textbf{Why the board demonstration needed a flip-flop.} At 50\,MHz, \code{data_ready} is high for
20\,ns. No LED can show that, so the demonstration adds one flip-flop that is set by the pulse and
cleared by reset, turning the event into a level a human can see.

\textbf{Whose problem that is.} Chapter~\ref{c3:ch}'s edge detector, run backwards. There the
problem was turning a level into a one-cycle event; here it is turning a one-cycle event back into a
level. Pulse against level is the same distinction both times, and it is why the extra flip-flop is
not a hack. \worth{1}

% ==========================================================================================
\answersection{7}{The counter that watches for a number}{10}

\answerpart{a}{3}
\textbf{Timer against counter.} A counter answers ``what is the current count?'' A timer answers
``has a certain amount of time passed?'' It is a counter plus two things: a \textbf{comparison}
against a fixed target, and a \textbf{clear} that makes it repeat. \worth{1}

\textbf{As gates.} One \textbf{XNOR per bit}, each comparing a counter bit against the corresponding
bit of the target: an XNOR outputs \code{1} when its inputs are the same, so each one is a single-bit
equality test. Their outputs are combined by a single \textbf{AND}, whose output is \code{timeout}.

\begin{booktable}{@{}p{6.5em}L@{}}
  \thead{Target width} & \thead{Gates} \\
  \midrule
  4 bits & 4 XNOR + 1 four-input AND = \textbf{5} \\
  26 bits & 26 XNOR + an AND combining 26 signals = \textbf{27}, or 51 with the AND built as a tree
    of 25 two-input gates \\
\end{booktable}

\noindent Twenty-six bits is what counting to 50,000,000 needs, and 26 XNORs across 52 inputs is
the drawing the VHDL exists to avoid. \worth{1}

\textbf{What insists on all four.} The AND gate. Three matching bits is not a match, and watching
the four XNOR outputs as the count runs makes that concrete: three of them are high on plenty of
counts that are not \code{1010}. The comparison is the AND, not the XNORs. \worth{1}

\answerpart{b}{3}
\textbf{With the clear.} The counter runs \code{0} through \code{10} inclusive and is then forced
back to \code{0}, so a full cycle is \code{11} edges: \code{10 - 0 + 1 = 11}.

\textbf{Without it.} The comparator still fires when the count reads \code{1010}, but nothing resets
the counter, so it carries on \code{1011, 1100, 1101, 1110, 1111} and wraps to \code{0000} on its
own --- four bits, sixteen states --- and does not read \code{1010} again until sixteen edges have
passed. So \code{timeout} fires every \textbf{16} edges. \credit{2 marks for both numbers with the
arithmetic}

\textbf{Which was asked for: 11.} Both are periodic; only one has the period you specified. And on
a 4-bit counter the two are close enough to look like a rounding error, which is exactly why this is
worth checking rather than eyeballing.

\textbf{The off-by-one in general.} A counter running \code{0} through \code{TICK_COUNT} visits
\code{TICK_COUNT + 1} values, so the period is \code{TICK_COUNT + 1} clock cycles, and for an exact
count you write \code{TICK_COUNT = seconds * 50,000,000 - 1}. The round numbers the book quotes
leave an error of between 0.02 and 0.2\,ppm, far below the board oscillator's own tolerance.
\worth{1}

\answerpart{c}{2}
The CircuitVerse \code{timeout} is a \textbf{combinational decode} of the count, so it is high
during the cycle the counter reads \code{1010}, and between edges it can \textbf{glitch}: while the
counter bits skew on a transition, the XNOR-and-AND network can momentarily see a match that never
existed.

The VHDL \code{timeout} is \textbf{registered} inside the clocked process, so it is high one cycle
later, during the cycle \code{counter} reads \code{0}, and it cannot glitch at all: it comes
straight off a flip-flop.

Same period, different phase, different glitch behaviour. \worth{1}

The distinction is exactly the one Chapter~\ref{c8:ch}'s \textbf{Moore-versus-Mealy} discussion
turns on: an output decoded combinationally from state bits can glitch while they skew, and what
makes an output glitch-free is registering it, the way \code{fsm_led}'s \code{LED_PROCESS} does.
``Moore'' on its own does not mean glitch-free. \worth{1}

\answerpart{d}{2}
\textbf{The rule broken: one clock, and everything else is an enable.} \code{timeout} is a signal
produced by logic, not a clock, and routing it into a clock port creates a second clock domain out
of it. \worth{1}

\textbf{What it costs.} A design with one clock is a design whose timing the tool can analyze; a
design with a logic-derived clock is not, so no meaningful Fmax exists for it. Everything downstream
is now asynchronous to everything else in the design, so every signal crossing back needs
Chapter~\ref{c4:ch}'s synchronizer machinery. On an FPGA it also either burns a global clock
resource or, worse, routes a clock over ordinary logic routing with skew nobody bounded. The book
names this the single most common mistake a programmer makes writing a baud generator.

\textbf{What to do instead.} Keep the one 50\,MHz clock everywhere and use \code{timeout} as a
one-cycle \textbf{enable} on a flip-flop that runs on it --- which is what \code{walking_led} does
with the same pulse, and what \code{fsm_led} does in \code{STATE_BLINK}. \worth{1}

% ==========================================================================================
\answersection{8}{A machine designed by hand, then written down}{13}

\answerpart{a}{5}
\textbf{State diagram.} Four states in a ring with fall-back arrows; \code{y = 1} in \code{S3} only,
and marked on the bubble rather than on any arrow, which is what makes it a Moore machine.

Every state has exactly two arrows out of it, one per value of \code{din}:

\begin{booktable}{@{}p{5.2em}LL@{}}
  \thead{From} & \thead{\code{din = 0}} & \thead{\code{din = 1}} \\
  \midrule
  \code{S0} (y=0) & \code{S0}, self-loop: a \code{0} matches nothing & \code{S1}: the first \code{1}
    \\
  \code{S1} (y=0) & \code{S0}: \code{10} is not a prefix of \code{110} & \code{S2}: \code{11} seen \\
  \code{S2} (y=0) & \code{S3}: \code{110} complete & \code{S2}, self-loop: still \code{11} \\
  \code{S3} (\textbf{y=1}) & \code{S0}: fresh start & \code{S1}: fresh start, on this \code{1} \\
\end{booktable}

\noindent Drawn, that is a ring \code{S0 -> S1 -> S2 -> S3 -> S0} for the advancing path, a
self-loop on \code{S0} and on \code{S2}, a fall-back arrow from \code{S1} to \code{S0}, and a
re-entry arrow from \code{S3} to \code{S1} on \code{din = 1}. That is eight arrows, two out of each
of the four states. \code{y = 1} is written \textbf{on the \code{S3} bubble}, not on any arrow, and
that is what makes it a Moore machine.

\credit{2 marks: 1 for two arrows out of every state, 1 for \code{y} marked on \code{S3} alone}

\textbf{State table}, \code{S0 = 00}, \code{S1 = 01}, \code{S2 = 11}, \code{S3 = 10}:

\begin{booktable}{@{}cccccc@{\hspace{1.5em}}L@{}}
  \thead{Q1} & \thead{Q2} & \thead{din} & \thead{Q1+} & \thead{Q2+} & \thead{y} &
    \thead{Meaning} \\
  \midrule
  0 & 0 & 0 & 0 & 0 & 0 & \code{S0} stays: a \code{0} matches nothing \\
  0 & 0 & 1 & 0 & 1 & 0 & \code{S0} to \code{S1}: saw the first \code{1} \\
  0 & 1 & 0 & 0 & 0 & 0 & \code{S1} back to \code{S0}: \code{10} is not a prefix \\
  0 & 1 & 1 & 1 & 1 & 0 & \code{S1} to \code{S2}: saw \code{11} \\
  1 & 1 & 0 & 1 & 0 & 0 & \code{S2} to \code{S3}: \code{110} complete \\
  1 & 1 & 1 & 1 & 1 & 0 & \code{S2} stays \\
  1 & 0 & 0 & 0 & 0 & \textbf{1} & \code{S3} to \code{S0}, fresh start \\
  1 & 0 & 1 & 0 & 1 & \textbf{1} & \code{S3} to \code{S1}, fresh start on this \code{1} \\
\end{booktable}

\noindent\textit{(2 marks for all eight rows; 1 if the four \code{S0}/\code{S1} rows are right and
the fall-backs are not)}

\textbf{\code{S2} on \code{din = 1}: stay in \code{S2}.} The last two bits seen are still
\code{1, 1}, which is exactly what \code{S2} means, so there is nowhere better to go: an input of
\code{1, 1, 1, 0} must still report a match. The fall-back and self-loop arrows are where these
machines go wrong, not the forward ones. \worth{1}

\answerpart{b}{3}
Maps with \mbox{\code{{Q1, Q2}}} down the rows in Gray-code order and \code{din} across:

\needspace{9\baselineskip}\noindent\textbf{\code{Q1+}}

\begin{narrowtable}{@{}ccc@{}}
  \thead{\code{Q1Q2} \textbackslash{} \code{din}} & \thead{\code{0}} & \thead{\code{1}} \\
  \midrule
  \code{00} &            &            \\
  \code{01} &            & \textbf{1} \\
  \code{11} & \textbf{1} & \textbf{1} \\
  \code{10} &            &            \\
\end{narrowtable}

\noindent Two groups of two: the whole \code{11} row gives \code{Q1Q2}, and the \code{din = 1}
column across the adjacent rows \code{01} and \code{11} gives \code{Q2·din}.

\begin{textdrawing}
Q1+ = Q1Q2 + Q2·din  =  Q2(Q1 + din)
\end{textdrawing}

\needspace{9\baselineskip}\noindent\textbf{\code{Q2+}}

\begin{narrowtable}{@{}ccc@{}}
  \thead{\code{Q1Q2} \textbackslash{} \code{din}} & \thead{\code{0}} & \thead{\code{1}} \\
  \midrule
  \code{00} &  & \textbf{1} \\
  \code{01} &  & \textbf{1} \\
  \code{11} &  & \textbf{1} \\
  \code{10} &  & \textbf{1} \\
\end{narrowtable}

\noindent One group of four: the entire \code{din = 1} column, in which only \code{din} is
constant.

\begin{console}
Q2+ = din
\end{console}

\textbf{\code{y}}, which reads only the state, confirming this is a valid Moore machine:

\begin{console}
y = Q1Q2'
\end{console}

\noindent\textit{(2 marks for the three equations, of which 1 is for \code{Q2+ = din} being spotted
as a group of four rather than written out as four terms)}

\textbf{The gates.} \code{Q1+} is one OR and one AND in the factored form, or two ANDs and an OR
unfactored. \code{Q2+} is a \textbf{wire}: \code{din} goes straight to the second flip-flop's
\code{D} input, with no gate at all. \code{y} is one AND with one input inverted. Two D flip-flops
hold \mbox{\code{{Q1, Q2}}}, and that is the whole machine: two flip-flops and three or four gates.
\worth{1}

\answerpart{c}{3}
\begin{vhdlcode}
library ieee;
use ieee.std_logic_1164.all;

entity seq_detect_110 is
    port(clock, reset_n: in  std_logic;
         din           : in  std_logic;
         y             : out std_logic);
end entity;

architecture behaviour of seq_detect_110 is
type state_t is (S0, S1, S2, S3);
signal reset_s2_n: std_logic;
signal state: state_t;
begin
    -- Moore output: reads the state and nothing else.
    y <= '1' when state = S3 else '0';

    reset_sync1: entity work.reset_sync
        port map(clock, reset_n, reset_s2_n);

    STATE_PROCESS: process(clock, reset_s2_n) is
    begin
        if (reset_s2_n = '0') then
            state <= S0;
        elsif (rising_edge(clock)) then
            case (state) is
                when S0 =>
                    if (din = '1') then state <= S1; else state <= S0; end if;
                when S1 =>
                    if (din = '1') then state <= S2; else state <= S0; end if;
                when S2 =>
                    if (din = '1') then state <= S2; else state <= S3; end if;
                when S3 =>
                    if (din = '1') then state <= S1; else state <= S0; end if;
                when others =>
                    state <= S0;
            end case;
        end if;
    end process;
end architecture;
\end{vhdlcode}

\noindent\textit{(2 marks: 1 for an enumerated \code{state_t} with a \code{case} covering every
state and a \code{when others}, 1 for the transitions agreeing with the candidate's own state
table)}

\code{reset_n} goes through \code{reset_sync} because this is a \textbf{top level} and owns the pin.
\code{din} does \textbf{not} get a synchronizer: it is synchronous serial data, already clocked by
this clock, and synchronizing it would only delay it and shift the sequence being detected. This is
the same split \code{seq_detect_mealy} makes.

\textbf{What the VHDL does not contain.} The \textbf{encoding} (\code{S0 = 00}, \code{S1 = 01},
\dots) and the \textbf{next-state equations} derived in (b). Both are gone. The synthesis tool
decides the encoding, and derives the gate-level logic from the \code{case} itself; it will
re-encode freely, often to one-hot, unless told otherwise. That is exactly the trade
Chapter~\ref{c8:ch} is about: you designed it by hand once so that you know what the tool is doing,
and then hand the tool the decision. \worth{1}

Credit an answer that also notes \code{y <= '1' when state = S3 else '0';} is a Moore output that
is still combinationally decoded, so it can glitch briefly while the state bits skew; registering it
inside a second clocked process is what makes it glitch-free.

\answerpart{d}{2}
\code{to_next_state} is now high on \textbf{every clock cycle the button is held}, not once per
press, and the \code{case} advances the state on every one of them.

\textbf{The arithmetic.} At 50\,MHz a press is 20\,ns per cycle. A brisk 100\,ms press is
\code{0.1 / 20e-9 = 5,000,000} clock cycles, so the machine takes 5,000,000 transitions around a
three-state ring. \code{5,000,000 mod 3 = 2}, so \emph{that} press lands two states on --- but
nobody can hold a button for exactly 100\,ms, and one cycle either way changes the answer. The state
it lands in is decided by how long you happened to hold the button, measured to 20\,ns, which from
the user's point of view is a coin toss. \worth{1}

\textbf{The fix.} Edge-detect the synchronized level, which is what \code{button_sync}'s third
flip-flop is for: compare the current synchronized value against the previous one, giving a pulse
high for exactly one clock cycle per press. That is the \code{X} of the hand-designed machine and the
\code{button_edge_s2} of \code{fsm_led}, and it is why the machine's input has to be an event rather
than a level. \worth{1}
